\documentclass[../../../../main.tex]{subfiles}
\begin{document}

図の長方形 $ABP_1P_5$ はある国境の町をあらわし，
各線分は道路をあらわす．
図の地点 $P_1$，$P_2$，$\cdots$，$P_9$ には外国への通路が開かれている．
いま，ある犯人が $B$ から外国に向って逃走しようとしているが，
この犯人は $P_j$ $(1 \leqq j \leqq 9)$ 以外の各交差点($B$を含む)において，
確率 $\displaystyle\frac{1}{2}$ ずつで真東または北東に通路をえらぶ．
この犯人を捕えるために3人の警官を $P_j$ $(1 \leqq j \leqq 9)$ のうちの
適当な3地点に配置しようとする．
どの3点に配置すれば，犯人を捕える確率 $p$ が最大となるか．
また，そのときの $p$ の最大値を小数第2位まで求めよ．
ただし，犯人は警官に出会わないで国境の地点に達すれば，
無事に逃げおおせるものとする．

\begin{figure}[htbp]
  \centering
  \begin{tikzpicture}[scale=0.8, >=stealth]
    \foreach \x in {0,...,8} {
      \draw[thick] (\x, 0) -- (\x, 4);
    }
    \foreach \y in {0,...,4} {
      \draw[thick] (0, \y) -- (8, \y);
    }
    \foreach \x in {0,...,7} {
      \foreach \y in {0,...,3} {
        \draw[thick] (\x, \y) -- (\x+1, \y+1);
      }
    }

    \node[left] at (0, 4) {$A$};
    \node[left] at (0, 0) {$B$};

    \node[below right] at (8, 0) {$P_1$};
    \node[right] at (8, 1) {$P_2$};
    \node[right] at (8, 2) {$P_3$};
    \node[right] at (8, 3) {$P_4$};
    \node[above right] at (8, 4) {$P_5$};
    \node[above] at (7, 4) {$P_6$};
    \node[above] at (6, 4) {$P_7$};
    \node[above] at (5, 4) {$P_8$};
    \node[above] at (4, 4) {$P_9$};

    \begin{scope}[shift={(1,-2)}]
      \draw[->, thick] (0, -0.8) -- (0, 0.8) node[above] {北};
      \draw[->, thick] (-0.8, 0) -- (0.8, 0) node[right] {東};
    \end{scope}
  \end{tikzpicture}
\end{figure}

\end{document}
